\documentclass[11pt, a4paper]{article}

\usepackage[utf8]{inputenc}
\usepackage[T1]{fontenc}
\usepackage[charter]{mathdesign} % Elegant, readable font
\usepackage{geometry}
\geometry{margin=1in}
\usepackage{enumitem}
\usepackage{amsmath}
\usepackage{booktabs}
\usepackage{hyperref}

\makeatletter
\renewcommand{\maketitle}{
  \begin{flushleft}
    {\huge \textbf{\@title} \par}
    \vspace{0.5cm}
    {\large \@author \par}
    \vspace{0.2cm}
    {\@date \par}
    \vspace{1cm}
    \hrule
  \end{flushleft}
}
\makeatother


\title{\textbf{Project Selection and Simple Corpus Analysis} \\ 
\large University of Ljubljana, EMAI, Spring Semester}
\author{Plabon Shaha, Pietro Sestito, Guillem Masdemont}
\date{\today}

\begin{document}

\maketitle

\section{Project Overview}
This report outlines our group selection and the initial stages of our corpus analysis project. Our focus is on character-level feature extraction for text classification, inspired by the architectures proposed by Zhang, Zhao, and LeCun (2015).

\section{Character-Level CNN Architecture}
The core approach involves treating text as a raw signal, processed through the following technical specifications:

\begin{itemize}[leftmargin=1.5in, labelsep=0.5in]
    \item[\textbf{Quantization}] One-hot encoding using an alphabet of $m=70$ characters (alphanumeric + punctuation).
    \item[\textbf{Activation}] ReLU activation: $h(x) = \max(0, x)$.
    \item[\textbf{Pooling}] Max-pooling layers to capture the most salient features across the sequence.
    \item[\textbf{Optimization}] SGD with momentum 0.9, starting $\eta=0.01$ (halved every 3 epochs).
    \item[\textbf{Minibatch}] 128 samples per iteration.
\end{itemize}

\section{Proposed Project Components}
\begin{enumerate}
    \item \textbf{Introduction:} Defining the scope of character-level vs. word-level modeling.
    \item \textbf{Related Work:} Reviewing Zhang et al. (2015) and existing CNN/RNN baselines.
    \item \textbf{Initial Ideas:} Exploring the impact of character quantization on noisy datasets.
    \item \textbf{Dataset/Corpus:} [Insert specific dataset name here, e.g., AG News or Yelp Reviews].
    \item \textbf{Repository:} A well-organized GitHub structure for reproducibility.
\end{enumerate}

\section{Related Work Summary}
> Zhang, Xiang, Junbo Zhao, and Yann LeCun. "Character-level convolutional networks for text classification." \textit{Advances in Neural Information Processing Systems}, 2015.

\end{document}